\documentclass[11pt]{article}
\usepackage[a4paper,margin=2.5cm]{geometry}
\usepackage{amsmath,amssymb}
\usepackage{algorithm}
\usepackage{algpseudocode}

\newcommand{\N}{\mathcal{N}}
\newcommand{\median}{\operatorname{median}}

\begin{document}

\section*{Supplementary Note: The KNN normalization algorithm}

KNN normalization estimates one size factor per cell for the protein data. The
key idea is that a cell and its neighbors in a KNN graph have, by construction,
a similar biological state, so the \emph{median} (over proteins) of the ratio
between a neighbor's and the cell's expression is a robust estimate of the
relative difference in sequencing depth between the two. Averaging this over
all neighbors gives the size factor of the cell, and the procedure is iterated
until the size factors converge.

\paragraph{Notation.}
Let $Y \in \mathbb{N}^{C \times P}$ be the raw protein count matrix of $C$ cells
and $P$ proteins, and $y_{ci}$ its entries. The KNN graph defines for every cell
$c$ a set of neighbors $\N(c)$ with $|\N(c)| = k_c$. Let $\pi$ be the
pseudocount.

\begin{algorithm}[h]
\caption{KNN normalization (\texttt{knn\_normalize})}
\begin{algorithmic}[1]
\Require Raw protein counts $Y$ ($C \times P$); optionally RNA counts $Z$
  ($C \times G$); source of the neighbors $\in \{\text{prot}, \text{rna},
  \text{existing}\}$; number of neighbors $k$; pseudocount $\pi$; maximum
  number of iterations $T_{\max}$; convergence threshold $\varepsilon$;
  log-transform flag.
\Ensure Normalized protein data $\tilde{X}$ ($C \times P$), and optionally the
  size factors of every iteration.

\Statex
\Statex \textbf{Step 1: KNN graph}
\If{$k$ is not given}
  \State $k \gets \max\bigl(15,\ \min(\operatorname{round}(C/20),\ 300)\bigr)$
\EndIf
\If{source = existing}
  \State Use the stored graph of connectivities.
\Else
  \If{source = prot}
    \State $A \gets Y$, \textbf{if} log-transform \textbf{then} $A \gets \log(1 + A)$
    \State Metric $\gets$ cosine; apply PCA \textbf{only if} $P \ge P_{\text{PCA}}$
  \Else
    \State $A \gets \log\bigl(1 + \text{library-size-normalize}(Z)\bigr)$
    \State Metric $\gets$ Euclidean; apply PCA
  \EndIf
  \State Build the $k$-nearest-neighbor graph of the cells from $A$ (or from its
    principal components) and keep its connectivities.
\EndIf
\State From the graph, form the list $\mathcal{E}$ of all (cell, neighbor) pairs
  $(c, n)$ with $n \in \N(c)$, and the number of neighbors $k_c$ of every cell.

\Statex
\Statex \textbf{Step 2: Initialization}
\State $x_{ci} \gets y_{ci} + \pi$ for all $c, i$
\State $S_{\text{tot}} \gets \sum_{c,i} x_{ci}$, \quad
  $\rho_c \gets \sum_{i} x_{ci}$
\State $f_c \gets 1$ for all $c$ \Comment{cumulative size factor}
\State $R_{cn} \gets \median_{i=1,\dots,P}\ \dfrac{x_{ni}}{x_{ci}}$
  for all $(c,n) \in \mathcal{E}$
  \Comment{computed only once}

\Statex
\Statex \textbf{Step 3: Iterative estimation of the size factors}
\For{$t = 1, \dots, T_{\max}$}
  \For{each cell $c$}
    \If{$k_c > 0$}
      \State $s_c \gets \dfrac{1}{k_c} \sum_{n \in \N(c)} R_{cn}$
    \Else
      \State $s_c \gets 1$
    \EndIf
  \EndFor
  \State $s_c \gets s_c \cdot \dfrac{S_{\text{tot}}}{\sum_{c'} s_{c'}\, \rho_{c'}\, f_{c'}}$
    for all $c$
    \Comment{keep the total counts constant}
  \State $s^{(t)} \gets s$; \quad $f_c \gets f_c \cdot s_c$ for all $c$
  \If{$t > 1$ \textbf{and} $\max_c \bigl| s^{(t)}_c - s^{(t-1)}_c \bigr| < \varepsilon$}
    \State \textbf{break}
  \EndIf
  \State $R_{cn} \gets R_{cn} \cdot \dfrac{s_n}{s_c}$ for all $(c,n) \in
    \mathcal{E}$ \Comment{update ratios without recomputing medians}
\EndFor

\Statex
\Statex \textbf{Step 4: Output}
\State $\tilde{x}_{ci} \gets f_c \cdot x_{ci}$ for all $c, i$
\If{log-transform}
  \State $\tilde{x}_{ci} \gets \log \tilde{x}_{ci}$
\EndIf
\State \Return $\tilde{X}$ (and, if requested, $f_c$, $s^{(1)}, \dots, s^{(t)}$ and the graph)
\end{algorithmic}
\end{algorithm}

\paragraph{Details and default parameters.}
\begin{itemize}
  \item \textbf{Number of neighbors.} The default $k = \max(15, \min(\operatorname{round}(C/20), 300))$
    scales with the number of cells, so that every cell has enough neighbors
    for a stable average, while the neighborhood stays local in large data sets.
  \item \textbf{Pseudocount.} The default is $\pi = 5$. It avoids divisions by
    zero in $x_{ni}/x_{ci}$. If the data contains no zeros, $\pi = 0$ can be used.
  \item \textbf{Median over proteins.} Using the median instead of the mean
    makes the ratio $R_{cn}$ robust to the few proteins that truly differ
    between a cell and its neighbor (biological variation), so that the ratio
    reflects the difference in depth.
  \item \textbf{Efficient update of the ratios.} After the size factors $s$ of an
    iteration have been applied, the expression of cell $c$ becomes $s_c x_{ci}$
    and that of its neighbor $s_n x_{ni}$. Since $s_c$ and $s_n$ are constant over
    proteins, the new median ratio is exactly $R_{cn}\, s_n / s_c$. Therefore the
    expensive median over proteins is computed only once (Step 2) and every
    iteration only needs $O(|\mathcal{E}|)$ operations. All size factors are
    applied to the data at once at the end, through the cumulative factor $f_c$.
  \item \textbf{Rescaling.} After each iteration the size factors are rescaled so
    that the total of the counts, $\sum_c s_c \rho_c f_c$, equals the original
    total $S_{\text{tot}}$. This fixes the arbitrary global scale of the size
    factors, which the median ratios alone do not determine.
  \item \textbf{Convergence.} The iterations stop when the maximum absolute
    change of any size factor between two consecutive iterations is below
    $\varepsilon = 0.0005$, or after $T_{\max} = 25$ iterations.
  \item \textbf{Cells without neighbors} receive the neutral size factor
    $s_c = 1$.
\end{itemize}

\end{document}
